\chapter{\IfLanguageName{dutch}{Stand van zaken}{State of the art}}%
\label{ch:stand-van-zaken}

% Tip: Begin elk hoofdstuk met een paragraaf inleiding die beschrijft hoe
% dit hoofdstuk past binnen het geheel van de bachelorproef. Geef in het
% bijzonder aan wat de link is met het vorige en volgende hoofdstuk.

% Pas na deze inleidende paragraaf komt de eerste sectiehoofding.

De beveiligingsrisico's van MCP-integraties in de financiële sector zijn pas goed te begrijpen tegen de achtergrond van hoe het protocol werkt en welke aanvalsvectoren het introduceert. Dit hoofdstuk brengt die context in kaart op basis van recente literatuur, en werkt toe naar de onderzoeksnood die de grondslag vormt voor de methodologie in Hoofdstuk~\ref{ch:methodologie}.

\section{Opkomst en risico's van MCP in Fintech}

De integratie van AI-agents in de financiële sector, voor toepassingen variërend van klantenservice tot transactieverwerking, leunt steeds meer op gestandaardiseerde protocollen. Het Model Context Protocol (MCP) verbindt AI-modellen met de bijbehorende bedrijfsapplicaties en databronnen. Hoewel deze standaardisatie integratie vergemakkelijkt, hebben \textcite{song2025protocolunveilingattackvectors} aangetoond dat het ecosysteem hierdoor nieuwe aanvalsvectoren introduceert die verder reiken dan het protocol zelf.

Die complexiteit brengt specifieke beveiligingsrisico's mee. Uit een systematische analyse door \textcite{guo2025systematicanalysismcpsecurity} blijkt dat MCP-implementaties vatbaar zijn voor diverse fundamentele kwetsbaarheden, waaronder \emph{prompt injection} via tool-beschrijvingen en \emph{privilege escalation}. Kwaadwillenden kunnen die zwakheden misbruiken. Zo beschrijven \textcite{Cato2025} in hun onderzoek naar "Living off AI"-aanvallen hoe MCP-servers gemanipuleerd kunnen worden omdat ze onvoldoende onderscheid maken tussen legitieme gebruikers en externe actoren, ook wel bekend als de \emph{confused deputy} aanval.

\section{Classificatie van kwetsbaarheden}

\textcite{OWASP2025} brachten de meest kritieke risico's in kaart in hun \emph{OWASP Top 10 for MCP}, met onder andere \emph{tool poisoning} en het gebrek aan degelijke audit logging als prominente bevindingen. Het is belangrijk om deze lijst te onderscheiden van de \emph{OWASP Top 10 for LLM Applications}~\autocite{OWASP2025}: de LLM-lijst richt zich op kwetsbaarheden in taalmodellen zelf (zoals onveilige output handling en training data poisoning), terwijl de MCP-lijst de beveiligingsrisico's van het protocol en de tool-interface beschrijft. Beide lijsten overlappen bij prompt injection en output sanitisatie, maar de MCP-specifieke risico's (zoals tool poisoning via kwaadaardige tool-beschrijvingen en het ontbreken van autorisatiechecks op tool-niveau) vallen buiten de scope van de LLM-lijst. In dit onderzoek hanteren we de \emph{OWASP Top 10 for MCP} als primaire taxonomie; waar relevant verwijzen we naar de LLM-lijst voor gedeelde risico's.

Voor fintech-bedrijven, die met gevoelige persoons- en transactiegegevens werken, vormen de MCP-specifieke kwetsbaarheden een direct risico op datalekken en compliance-schendingen. De OWASP MCP-taxonomie categoriseert aanvalsvectoren die algemene frameworks zoals de OWASP Top 10 voor webapplicaties niet dekken.

\textcite{hou2025mcplandscape} bieden een aanvullende taxonomie die MCP-dreigingen indeelt in vier aanvalstypen (kwaadwillende ontwikkelaars, externe aanvallers, misbruikende gebruikers en structurele implementatiefouten), verdeeld over 16 concrete bedreigingsscenario's. Het framework in dit onderzoek richt zich primair op de eerste en vierde categorie: kwetsbaarheden die een ontwikkelaar onbewust introduceert in tool-beschrijvingen en inputschema's. \textcite{huang2026mcpthreatmodeling} tonen via STRIDE-dreigings\-modellering aan dat de meeste MCP-clients geen statische validatie uitvoeren op ontvangen tool-metadata, wat precies de lacune is die de SAST-module in dit onderzoek opvult.

\section{Status Quo: Beperkingen van handmatige validatie}

Ondanks de ernst van de dreigingen, ontbreken er momenteel gestandaardiseerde methodologieën voor validatie. De huidige praktijk in het bedrijfsleven leunt sterk op handmatige en herhaalde code reviews en het gebruik van ad-hoc tools. \textcite{Uproot2025} bieden weliswaar handleidingen voor het uitvoeren van MCP-pentests, maar deze processen blijven arbeidsintensief en foutgevoelig. Een handmatige security review van een MCP-server implementatie kan gemiddeld 8 tot 16 uur in beslag nemen per iteratie, afhankelijk van de complexiteit van de server \autocite{Uproot2025}. In teams die frequent deployen, maakt dat grondige beveiligingscontroles moeilijk vol te houden.

Bestaande tools bieden slechts gedeeltelijke oplossingen. Hoewel er open-source initiatieven zijn, zoals de scanning tool van \textcite{Invariant2025}, ontbreekt vaak een integratie in een breder security framework dat geschikt is voor enterprise-omgevingen. Bedrijven vallen hierdoor terug op algemene tools zoals Burp Suite of Semgrep. Burp Suite kan JSON-RPC-verkeer onderscheppen maar mist de contextuele kennis om kwaadaardige tool-beschrijvingen te herkennen. Semgrep detecteert codepatronen maar kan de runtime-interactie tussen een agent en een MCP-server niet analyseren. Beide tools vullen daarmee elk een ander deel van het aanvalsvlak in, maar niet het MCP-specifieke deel.

Het ontbreken van zo'n framework is een onderbelicht gebied in zowel de literatuur als de praktijk, al evolueert het domein snel: de meeste relevante publicaties dateren uit 2025 en het MCP-ecosysteem is te jong voor definitieve generalisaties.

\section{Bestaande tools voor MCP-beveiligingstesting}

Er bestaan al enkele MCP-specifieke tools die gedeeltelijke dekking bieden. De meest directe is MCPSafetyScanner, ontwikkeld door \textcite{radosevich2025mcpsafetyaudit}. MCPSafetyScanner hanteert een agentische benadering: een LLM-agent genereert adversariale testscenario's en evalueert de respons van de doelserver. Dit maakt de tool contextbewust en flexibel in de keuze van testcases, maar ook niet-deterministisch: twee opeenvolgende runs op dezelfde server kunnen verschillende bevindingen opleveren, afhankelijk van het gedrag van het onderliggende model. Bovendien vereist de tool een actieve API-verbinding met een extern taalmodel, wat in gereguleerde fintech-omgevingen toestemmings- en dataprivacyvraagstukken meebrengt.

Het framework in dit onderzoek kiest bewust voor een deterministische aanpak. Patroonherkenning op basis van vaste detectieregels en vooraf gedefinieerde fuzzingpayloads vervangen de LLM-redenering van MCPSafetyScanner. Dit gaat ten koste van semantische flexibiliteit, maar garandeert reproduceerbaarheid (NF-2) en volledige onafhankelijkheid van externe diensten.

Een complementaire benadering biedt MCPTox~\autocite{wang2025mcptox}: een benchmark van 45 echte MCP-servers en 353 authentieke tools waarmee de robuustheid van \emph{LLM-agenten} tegen tool poisoning wordt gemeten. MCPTox rapporteert een aanvalssuccesratio van 72{,}8\% op o1-mini, wat de ernst van tool poisoning in productieomgevingen empirisch aantoont. Het onderscheid in evaluatieperspectief is hier van belang: MCPTox test of een model bestand is tegen een aanval, terwijl dit framework test of een server de aanvalsvector überhaupt \emph{bevat}. Beide invalshoeken vullen elkaar aan.

\section{Geautomatiseerde beveiligingstesting als oplossing}

De softwarebeveiligingsliteratuur beschrijft een reeks geautomatiseerde testmethoden die de nadelen van handmatige validatie kunnen opvangen. \emph{Static Application Security Testing} (SAST) analyseert broncode zonder uitvoering, wat het geschikt maakt voor vroege detectie van onveilige patronen in tool-beschrijvingen en configuratiebestanden \autocite{chess2004static}. \emph{Dynamic Application Security Testing} (DAST) test een draaiende applicatie door gemanipuleerde invoer te sturen en de respons te observeren. Deze aanpak is direct toepasbaar op de JSON-RPC-interface van MCP-servers. \emph{Fuzzing} ten slotte genereert automatisch grote volumes onverwachte of grensgevallen-invoer om onvoorziene kwetsbaarheden bloot te leggen \autocite{manes2019art}.

Toegepast op MCP vullen deze drie methoden elkaar aan: SAST detecteert statische patronen in tool-metadata en configuratiebestanden, terwijl DAST en fuzzing runtime-gedrag blootleggen dat niet uit de broncode alleen af te leiden is. Zo kunnen gemanipuleerde tool-beschrijvingen pas tijdens uitvoering hun effect tonen op het gedrag van een agent. Een framework dat deze drie benaderingen combineert, dekt zowel de statische als de dynamische aanvalsvlakken van MCP-servers.

\section{Kritische noot bij de literatuur}

De literatuur over MCP-beveiliging is jong: de meeste aangehaalde publicaties dateren uit 2025, en de OWASP MCP Top 10 werd pas gepubliceerd nadat MCP zelf breed werd geadopteerd. Dat heeft twee gevolgen voor de interpretatie. Sommige claims in de literatuur (waaronder de classificaties van \textcite{hou2025mcplandscape} en de dreigingsmodellen van \textcite{huang2026mcpthreatmodeling}) zijn nog niet empirisch gevalideerd op grote schaal. Daarnaast evolueert het MCP-ecosysteem snel: nieuwe protocol-versies, uitbreidingen en adoptie door grote platformen kunnen de relevantie van specifieke aanvalsvectoren snel doen verschuiven. Uitspraken over "het risico van MCP" moeten dan ook gelezen worden als een momentopname van de situatie in 2025, niet als definitieve conclusies.


